% ============================================================
% 分片：s07  对应原文 PDF 第 72--76 页
% 负责章节：附录 B 收尾（式 B.26–B.27）；附录 C Landau–Ginzburg $B$-models and Matrix Factorizations：
%           C.1 $B$-type quantum mechanics；C.2 Boundary conditions in $B$-twisted Landau–Ginzburg models；
%           C.3 Matrix factorizations from $B$-branes
% 译者：translator-08   审校状态：初稿
% 注：附录模式已在 s04 生效，本分片的 \section 自动编号为 C。
% ============================================================
% ==== 原文 PDF 第 72 页 ====
% [跨页接续: 原文第71页/72页，本句承接 s06 末段（第 71 页末仅有 "In more detail, the index" 半句，按语义完整取于此）]
更详细地说，Dirac 算子 $L$ 的指标意味着，梯度流 $p\to q$ 的模空间（至少在一般超势 $h$ 下）的维数为 $\mu(q)-\mu(p)$。当前情形下该模空间是 1 维的。由于总可以平移时间参数，我们发现这类流的约化模空间，即完整模空间对这些平移的商，是 0 维的，即一些离散点。综合起来，我们得出结论：超对称基态空间可等同于 Morse–Smale–Witten 复形的上同调：
\begin{equation}\tag{B.26}
  \mathcal{H}_{MSW}=\bigoplus|p\rangle,\qquad \bar Q|p\rangle=\sum_{q\,|\,\mu(q)=\mu(p)+1}m_{qp}|q\rangle.
\end{equation}
% [Tutor Note: (B.26) 把微扰基态 $|p\rangle$ 按 Morse 指标 $\mu(p)$ 分次，微分 $\bar Q$ 仅连接指标相差 $1$ 的临界点，这正是 Morse 复形的链层定义；"取 $\bar Q$-上同调"即取该复形的同调，因此基态计数（Witten 指标）等于 $M$ 的 Euler 示性数。]

最后这一断言需要一些论证——至少应当说明 Morse–Smale–Witten 复形确实是一个复形，即为何 $\bar Q^2=0$。具体地，有
\begin{equation}\tag{B.27}
  \bar Q^2|p\rangle=\sum_{r\,|\,\mu(r)=\mu(q)+1}\sum_{q\,|\,\mu(q)=\mu(p)+1}m_{rq}m_{qp}|r\rangle,
\end{equation}
% [Tutor Note: 两个求和是嵌套的：外层固定 $r$，内层对满足 $\mu(q)=\mu(p)+1$ 的中间临界点求和；注意外层条件写成 $\mu(r)=\mu(q)+1$，它与 $\mu(q)=\mu(p)+1$ 联立后给出 $\mu(r)=\mu(p)+2$，这正是非零条件。核对时可先用指标条件化简求和范围。]
于是必须检验：对每个满足 $\mu(r)=\mu(p)+2$ 的临界点 $r$，$\sum_q m_{rq}m_{qp}$ 为零，其中求和遍历满足 $\mu(q)=\mu(p)+1=\mu(r)-1$ 的中间临界点 $q$。该和式只有当存在临界点 $q$ 及梯度流 $\gamma_{qp}:p\to q$ 与 $\gamma_{rq}:q\to r$ 时才可能非零，故设此情形成立。如 [48, Section 10.6] 所述，复合路径 $\gamma_{rq}+\gamma_{qp}$ 可被 $p\to r$ 的梯度流任意好地逼近；此外，由于这类轨迹模空间的期望维数为 $\mu(r)-\mu(p)=2$，它属于这样的梯度流的一个 2 维族。一旦对时间参数的整体平移作商，我们便得到 1 维约化模空间——它是流 $p\to r$ 的完整模空间的某个连通分支。

再次地，指标理论意味着轨迹 $p\to q$ 的约化模空间的给定分支必然是光滑、无边界的 1 维流形。有两种可能：要么模空间紧致，从而是 $S^1$ 的一个拷贝；要么非紧，从而是 $\mathbb{R}$ 的一个拷贝。对当前分支，我们知道它不可能紧致——复合路径 $\gamma_{rq}+\gamma_{qp}$ 是该模空间的一个极限点。特别地，必然存在第二条极限轨迹（对应于 $\mathbb{R}$ 上的另一个极限点），它同样对应于复合轨迹 $\gamma_{rq'}+\gamma_{q'p}$，其中 $q'$ 是一个（可能不同的）不动点，满足 $\mu(q')=\mu(p)+1$。而且，$\gamma_{rq'}+\gamma_{q'p}$ 与 $\gamma_{rq}+\gamma_{qp}$ 必有相反的（相对）定向。因此，复合轨迹 $\gamma_{rq}+\gamma_{qp}$ 对 $\bar Q^2|p\rangle$ 的贡献必然被第二条复合轨迹 $\gamma_{rq'}+\gamma_{q'p}$ 抵消。
% [Tutor Note: 这段是 $\bar Q^2=0$ 的标准证明：轨迹模空间的维数为 $\mu(q)-\mu(p)=1$，商去时间平移（本身贡献 1 维）后得到 0 维点集，故带符号计数 $m_{qp}$ 有定义；而 1 维约化模空间的边界（流线在中间时刻断裂为两段）必由偶数个点组成，且相邻两点的相对定向相反，故两项贡献相消。物理上这对应"1 维瞬子模空间的边界由断裂瞬子组成"，与 B.2 节的局域化论证同一逻辑。]

\section{Landau–Ginzburg $B$ 模型与矩阵分解}
\begin{keytakeaway}
\textbf{核心物理图像：}$B$ 型扭变的 Landau–Ginzburg 模型的 1/2-BPS 边界条件（$B$ 膜）不能直接由体理论给出：体超势 $W$ 在边界上留下非零的全导数项，必须引入边界费米自由度将其吸收；把 $W$ 分解为 $E^a_{1d}J_{1d,a}$ 正是矩阵分解的物理来源。
\textbf{数学工具：}1d $\mathcal{N}=2$ 超空间与 $E$/J 项、$B$ 型扭变超荷 $\bar Q_B$ 的上同调、$\mathbb{Z}$-分次模与自同态、DG 范畴 $MF(W)$。
\textbf{主要结论：}$B$ 膜的范畴（其交界面上的局域算子为态射）等价于超势 $W$ 的矩阵分解的（DG-）范畴；矩阵分解 $(E,\delta_E)$ 满足 $\delta_E^2=W(\phi|)\mathrm{id}_E$，膜间局域算子即满足 $\bar Q_B$-闭条件的模映射。
\end{keytakeaway}

本附录回顾 $B$ 型扭变的 Landau–Ginzburg 模型的若干侧面，以及 Kontsevich 提出的、在边界条件语境中矩阵分解的来源 [78, Section 7]。我们先在 C.1 节回顾 1d $\mathcal{N}=2$ 的 $B$ 型量子力学的若干方面。在 C.2 节中，我们把 2d Landau–Ginzburg 模型改写为含 $B$ 型扭变超荷 $\bar Q_B=\bar Q_++\bar Q_-$\footnote{仅在本附录中，我们把 $B$ 型扭变超荷记作 $\bar Q_B$，以便比较我们对 $\mathcal{N}=2$ 量子力学的讨论与 $B$ 型扭变 2d $\mathcal{N}=(2,2)$ 的约化。} 的 1d $\mathcal{N}=2$ 代数，并通过引入对边界超势作分解的边界 Fermi 多重态来描述 $B$ 膜。最后，在 C.3 节中我们说明矩阵分解范畴如何作为 $B$ 膜交汇处的一个有用描述而出现。

% ==== 原文 PDF 第 73 页 ====
\subsection{$B$ 型量子力学}
我们如 B.2 节那样从坐标为 $t,\theta,\bar\theta$ 的超空间 $\mathbb{R}^{1|2}$ 出发；超对称生成元 $Q,\bar Q$ 与超导数 $D,\bar D$ 的约定与 B.1 节相同。我们将考虑的第一类不可约超场是手征超场 $\Phi$，它是玻色子的且满足 $\bar D\Phi=0$。手征超场有如下分量展开：
\begin{equation}\tag{C.1}
  \Phi^n=\phi^n+\theta\psi^n+\theta\bar\theta(-\partial_t\phi^n).
\end{equation}
若 $\Phi^n$（从而 $\phi^n$）具有 $R$ 荷 $r_n$，则 $\psi^n$ 具有 $R$ 荷 $r_n+1$。我们要考虑的第二类不可约超场是费米子的，称为 Fermi 超场 $\Gamma$。Fermi 多重态的最一般形式依赖于手征超场的一个全纯函数 $E^a(\phi)$ 的选取，该函数称为 $E$ 项，并满足一个修改后的手征性约束：$\bar D\Gamma^a-E^a(\Phi)=0$。该约束意味着 $\Gamma^a$ 可展开为
\begin{equation}\tag{C.2}
  \Gamma^a=\gamma^a+\theta g^a-\bar\theta(E^a(\phi))+\theta\bar\theta(-\partial_t\gamma^a+\psi^n\partial_nE^a(\phi)).
\end{equation}
其中 $\partial_n=\partial_{\phi^n}$。这样的 Fermi 多重态称为 $E$ 型 Fermi 多重态。为保持 $U(1)_R$ $R$ 对称性，我们假定 $E^a$ 具有齐次次数 $r_a+1$，其中 $r_a$ 是 $\Gamma^a$ 的 $R$ 荷。

由上可知，超荷 $Q$ 对这些超场的分量及其复共轭的作用如下：
\begin{equation}\tag{C.3}
  \begin{aligned}
  &Q\phi^n=\psi^n\quad Q\psi^n=0\quad Q\gamma^a=g^a\quad Qg^a=0\\
  &Q\bar\phi^{\bar n}=0\quad Q\bar\psi^{\bar n}=2\partial_t\bar\phi^{\bar n}\quad Q\bar\gamma^{\bar a}=-\bar E^{\bar a}\quad Q\bar g^{\bar a}=-2(\partial_t\bar\gamma^{\bar a}+\tfrac{1}{2}\bar\psi^{\bar n}\partial_{\bar n}\bar E^{\bar a})
  \end{aligned}
\end{equation}
其中 $\partial_{\bar n}=\partial_{\bar\phi^{\bar n}}$。类似地，$\bar Q$ 的作用为：
\begin{equation}\tag{C.4}
  \begin{aligned}
  &\bar Q\phi^n=0\quad \bar Q\psi^n=2\partial_t\phi^n\quad \bar Q\gamma^a=E^a\quad \bar Qg^a=2(\partial_t\gamma^a-\tfrac{1}{2}\psi^n\partial_nE^a)\\
  &\bar Q\bar\phi^{\bar n}=\bar\psi^{\bar n}\quad \bar Q\bar\psi^{\bar n}=0\quad \bar Q\bar\gamma^{\bar a}=-\bar g^{\bar a}\quad \bar Q\bar g^{\bar a}=0
  \end{aligned}
\end{equation}
% [Tutor Note: (C.3)、(C.4) 可直接由 (C.1)、(C.2) 的分量展开读出：$Q$ 对应 $\partial_\theta$、$\bar Q$ 对应 $\partial_{\bar\theta}$ 在 $\theta=\bar\theta=0$ 处的取值，而 $\partial_t$ 项来自 $\theta\bar\theta$ 分量中 $\theta,\bar\theta$ 的各阶导数。核对时建议把 $Q$ 作用在 (C.1) 上逐项验证，特别是 $Q\bar\psi^{\bar n}=2\partial_t\bar\phi^{\bar n}$ 中的因子 $2$ 来自 $\{\bar Q,Q\}=2H$（B.3 节约定）。]

我们可以把手征多重态与 Fermi 多重态的一个在 (C.3)、(C.4) 变换下不变的简单作用量写成超空间上的积分。每个 Fermi 多重态都有一个相伴的 $E$ 项，此外我们再引入第二个全纯函数 $J_a(\phi)$，即 $J$ 项。为保持 $U(1)_R$，需要 $J_a$ 是 $1-r_a$ 次齐次的。理论的作用量可写成超空间积分：
\begin{equation}\tag{C.5}
  \begin{aligned}
  S&=\frac{1}{4}\int dtd^2\theta\left(\delta_{n\bar n}(D\Phi^n)(\bar D\bar\Phi^{\bar n})-\delta_{a\bar a}\Gamma^a\bar\Gamma^{\bar a}\right)-\frac{1}{4}\int dtd\theta\,\Gamma^aJ_a(\Phi)\Big|_{\bar\theta=0}+\text{c.c.}\\
  &=\int dt\Big[|\partial_t\phi^n|^2-\frac{1}{2}\bar\psi_n\partial_t\psi^n-\frac{1}{2}\bar\gamma_a\partial_t\gamma^a+\frac{1}{4}E^a\bar E_a-\frac{1}{4}g^a\bar g_a-\frac{1}{4}g^aJ_a(\phi)-\frac{1}{4}\bar g_a\bar J^a(\bar\phi)\\
  &\qquad+\frac{1}{4}\left(\bar\gamma_a\psi^n\partial_nE^a+\gamma^a\psi^n\partial_nJ_a-\gamma^a\bar\psi_n\bar\partial^n\bar E_a-\bar\gamma_a\bar\psi_n\bar\partial^n\bar J^a\right)\Big],
  \end{aligned}
\end{equation}
其中我们用了度规 $\delta_{\bar n n}$ 与 $\delta_{\bar a a}$ 消去所有带杠指标，且 $\bar\partial^n=\delta^{n\bar n}\partial_{\bar\phi^{\bar n}}$。如同 2d $\mathcal{N}=(2,2)$ 的手征多重态，可以把动能项 $\delta_{\bar n n}(\bar D\bar\Phi^{\bar a})(D\Phi^a)$ 换成 $\bar DDK(\Phi,\bar\Phi)$，其中 $K(\phi,\bar\phi)$ 是更一般的凯勒势；等价地，可把 $\delta_{\bar n n}$ 换成更一般的凯勒度规 $G_{\bar n n}(\Phi,\bar\Phi)$；对 Fermi 多重态也可类似地引入非平凡的 Hermitian 度规 $H_{\bar a a}(\Phi,\bar\Phi)$。
% [Tutor Note: (C.5) 第二行中 $\frac{1}{4}E^a\bar E_a-\frac{1}{4}g^a\bar g_a$ 的组合来自消去辅助场之前的动能项；把 $g^a=-\bar J^a$、$\bar g_a=-J_a$ 代入后，$\frac{1}{4}(E^a\bar E_a-\bar J^aJ_a)$ 才是真正的势能，这一点在 C.2 节讨论 $\bar Q$-闭条件时会用到。]
% ==== 原文 PDF 第 74 页 ====
注意，来自 Fermi 多重态的玻色子 $g,\bar g$ 是辅助场，其运动方程将它们专设为 $g^a=-\bar J^a$ 与 $\bar g_a=-J_a$。积分掉这些辅助场后，我们看到作用量 $S$ 与超对称变换在同时交换费米子 $\gamma\leftrightarrow\bar\gamma$ 以及 $E$ 项与 $J$ 项 $E\leftrightarrow J$ 之下不变，这是 2d $\mathcal{N}=(0,2)$ 理论中费米子 $T$ 对偶的 1d 类比，参见 [45, Appendix A]。

为使 (C.5) 中的作用量超对称，即让所给出的 $Q,\bar Q$ 作用成为真正的对称性，当且仅当 $E^aJ_a$ 为常数。动能项不变的理由与 B.2 节相同，因此问题出在 $J$ 项。例如：
\begin{equation}\tag{C.6}
  \begin{aligned}
  &Q\int dtd\theta\,\Gamma^aJ_a(\Phi)\Big|_{\bar\theta=0}=0\\
  &\bar Q\int dtd\theta\,\Gamma^aJ_a(\Phi)\Big|_{\bar\theta=0}=\int dt\left(2\partial_t(\gamma^aJ_a)-\psi^n\partial_n(E^aJ_a)\right)
  \end{aligned}
\end{equation}
% [Tutor Note: (C.6) 第二式中的 $2\partial_t(\gamma^aJ_a)$ 是全导数项，不贡献物理；真正 obstruction 是 $-\psi^n\partial_n(E^aJ_a)$。因此要求 $\partial_n(E^aJ_a)=0$ 即 $E^aJ_a$ 为常数；若再要求 $U(1)_R$ 不变，则常数必须为零，得到著名的 $E^aJ_a=0$。这正是后面用边界费米多重态"因子化" $W$ 的动机。]
在 $Q$ 下的变分平凡地不变，但在 $\bar Q$ 下的变分几乎是全导数——只要我们要求 $\partial_n(E^aJ_a)=0$，即 $E^aJ_a$ 为常数，它就确实是全导数。若想保持 $U(1)_R$ $R$ 对称性，则该常数必须为零，于是得到著名关系 $E^aJ_a=0$。我们将看到，一旦把这些理论放到 2d 理论的边界上，这一条件就可以放宽：体作用量变分给出的边界项可用来抵消边界作用量的变分。

量子化上述经典理论相当直接。理论的 Hilbert 空间可等同于玻色子 $\phi^n,\bar\phi_n$ 以及一半费米子（例如 $\bar\psi_a$ 与 $\bar\gamma_n$）的多项式，$\mathcal{H}\simeq\mathbb{C}[\phi^n,\bar\phi_n,\bar\psi_a,\bar\gamma_n]$。\footnote{更确切地说，要得到平方可归一的态，我们应当为一个 $U(1)$ 味对称性引入实质量参数 $m$，参见 [150, Section 2]。这个实质量参数在本文所述的对 $\phi,\bar\phi$ 的多项式依赖之上诱导出 $\sim e^{|\phi|^2}$ 的指数压低。}超荷 $Q,\bar Q$ 表示为如下微分算子：
\begin{equation}\tag{C.7}
  Q=\partial_{\bar\psi_n}\partial_n-\bar J^a\bar\gamma_a-\bar E_a\partial_{\bar\gamma_a}\qquad \bar Q=\bar\psi_n\bar\partial^a+E^a\bar\gamma_a+J_a\partial_{\bar\gamma_a}.
\end{equation}
% [存疑: (C.7) 中第二式印作 $\bar Q=\bar\psi_n\bar\partial^a+\dots$，按指标应与 (C.6) 中的 $\psi^n\partial_n$ 配对而写成 $\bar\partial^n$；此处的上标 $a$ 疑为原文笔误，照原文排印]
$\mathcal{H}$ 的 $\bar Q$-上同调有一个自然的代数解释。第一项把 $\bar\phi_n$ 与 $\bar\psi_n$ 从上同调中移除，我们把剩余部分 $\mathbb{C}[\phi^n,\bar\gamma_a]$ 解释为一个 $\mathbb{Z}$-分次的 $\mathbb{C}[\phi^n]$-模 $\mathcal{E}$，其 $\mathbb{Z}$-分次来自 $U(1)_R$ $R$ 对称性。$\bar Q$ 中的第二项对应于 $\mathcal{E}$ 上的一个微分 $\delta$，于是 $\mathcal{H}$ 的 $\bar Q$-上同调可等同于这个 $\mathbb{C}[\phi]$-模复形的上同调：$H^\bullet(\mathcal{H})\simeq H^\bullet(\mathcal{E},\delta)$。\footnote{这也可以用复微分几何的语言表述如下。把费米子 $\bar\psi_n$ 等同于微分形式 $d\bar\phi_n$，于是 $\bar Q$ 的第一项就是靶空间上的 Dolbeault 微分 $\bar\partial=d\bar\phi_n\bar\partial^n$。接着把费米子 $\bar\gamma^{\bar n}$ 中的多项式等同于某个 $\mathbb{Z}$-分次向量丛 $E$ 的截面，而 $\bar Q$ 的其余部分把它变成一个向量丛复形 $\mathcal{E}$。综合起来，我们发现 $\mathcal{H}$ 的 $\bar Q$-上同调正是取值于 $\mathcal{E}$ 的第零个 Dolbeault 上同调群。}

\subsection{$B$ 型扭变 Landau–Ginzburg 模型中的边界条件}
描述 $B$ 型扭变 2d $\mathcal{N}=(2,2)$ 理论的（1/2-BPS）边界条件最有效的方式之一，是选取 2d $\mathcal{N}=(2,2)$ 超对称代数中一个含 $B$ 型扭变超荷 $\bar Q_B$ 的 1d $\mathcal{N}=2$ 子代数，然后把 2d 超空间作用量写成对「空间」$s$ 与 1d $\mathcal{N}=2$ 超空间的积分。对 $B$ 型扭变超荷 $\bar Q_B=\bar Q_++\bar Q_-$，共轭超荷是 $Q_B=Q_++Q_-$。它们共同生成所需的 1d $\mathcal{N}=2$ 代数 $\{Q_B,\bar Q_B\}=-2\partial_t$。

% ==== 原文 PDF 第 75 页 ====
相对于这个 1d $\mathcal{N}=2$ 子代数，一个 2d $\mathcal{N}=(2,2)$ 手征多重态 $\Phi^n$ 分解为一个手征多重态 $\Phi^n$ 与一个 $E$ 型 Fermi 多重态 $\Psi^n$，其中 $E^n_{2d}=2i\partial_s\phi^n$。更详细地说，我们定义 $\theta=\frac{1}{2}(\theta_++\theta_-)$ 与 $\eta=\frac{1}{2}(\theta_+-\theta_-)$，以及相应的导数。$\Phi^n$ 上的 2d 手征条件意味着 $\Phi^n_{2d}:=\Phi^n|_{\eta=\bar\eta=0}$ 满足 1d 手征性条件 $\bar D\Phi^n_{2d}=0$，其中 $\bar D=\bar D_++\bar D_-$。我们还发现 $\Gamma^n_{2d}:=(D'\Phi^n)|_{\eta=\bar\eta=0}$（其中 $D'=D_+-D_-$）正是上述 $E$ 型 Fermi 多重态：
\begin{equation}\tag{C.8}
  \bar D\Gamma^n_{2d}=i\bar DD'\Phi^n|_{\eta=\bar\eta=0}=i\{\bar D,D'\}\Phi^a|_{\eta=\bar\eta=0}=2i\partial_s\Phi^n_{2d}.
\end{equation}
% [存疑: (C.8) 中间的 $\Phi^a$ 按上下文应为 $\Phi^n$（左右两端均带指标 $n$，且 $\{\bar D,D'\}$ 作用在同一个多重态上）；照原文排印]
用分量场表示，我们有
\begin{equation}\tag{C.9}
  \Phi^n_{2d}=\phi^n+\theta\psi^n+\theta\bar\theta(-\partial_t\phi^n)\qquad \Gamma^n_{2d}=\gamma^n+\theta g^n-\bar\theta(2i\partial_s\phi^n)+\theta\bar\theta(-i\partial_t\gamma^n+2i\partial_s\psi^n)
\end{equation}
其中 $\psi^n=\psi^n_++\psi^n_-$，$\gamma^n=(\psi^n_+-\psi^n_-)$，且 $g^n=2F^n$。$Q_B,\bar Q_B$ 的作用可直接从 (C.3)、(C.4) 读出：
\begin{equation}\tag{C.10}
  \begin{aligned}
  &Q_B\phi^n=\psi^n\quad Q_B\psi^n=0\quad Q_B\gamma^n=g^n\quad Q_Bg^n=0\\
  &Q_B\bar\phi_n=0\quad Q_B\bar\psi_n=2\partial_t\bar\phi_n\quad Q_B\bar\gamma_n=2i\partial_s\bar\phi_n\quad Q_B\bar g_n=-2(\partial_t\bar\gamma_n-i\partial_s\bar\psi_n)
  \end{aligned}
\end{equation}
\begin{equation}\tag{C.11}
  \begin{aligned}
  &\bar Q_B\phi^n=0\quad \bar Q_B\psi^n=2\partial_t\phi^n\quad \bar Q_B\gamma^n=2i\partial_s\phi^n\quad \bar Q_Bg^n=2(\partial_t\gamma^n-i\partial_s\psi^n)\\
  &\bar Q_B\bar\phi_n=\bar\psi_n\quad \bar Q_B\bar\psi_n=0\quad \bar Q_B\bar\gamma_n=-\bar g_n\quad \bar Q_B\bar g_n=0
  \end{aligned}
\end{equation}
% [Tutor Note: (C.10)、(C.11) 中的 $\partial_s$ 项正是"2d 理论沿空间方向约化"的痕迹：$E^n_{2d}=2i\partial_s\phi^n$ 把 2d 的 $\partial_s$ 变为 1d 的 $E$ 项，而 $Q_B,\bar Q_B$ 与 2d 的 $Q_\pm$ 通过 $Q_B=Q_++Q_-$ 相联系。核对因子 $2i$ 时应回到 (C.8) 的 $\{\bar D,D'\}=2\partial_s$。]

通过对 $\eta$ 积分，超势项 $\int d^2\theta\,W(\Phi)$ 变成一个 $J$ 型超势，其中 $J_{2d,n}=\partial_nW(\phi)$。如上所述，我们假设 2d 手征多重态具有 $U(1)_R$ 荷 $r_a$，使得 $W$ 的 $R$ 荷为 2。于是 2d 作用量可表示为
\begin{equation}\tag{C.12}
  \begin{aligned}
  S_{2d}&=\frac{1}{4}\int dtd^2\theta\left(\int ds(D\Phi^n_{2d})(\bar D\bar\Phi_{2d,n})-\Gamma^n_{2d}\bar\Gamma_{2d,n}\right)\\
  &\qquad+\frac{1}{4}\int dtd\theta\left(\int ds\,\Gamma^n_{2d}J_{2d,n}\right)\Big|_{\bar\theta=0}+\text{c.c.}
  \end{aligned}
\end{equation}
注意，当没有空间边界时，我们得到 $\int ds\,E^n_{2d}J_{2d,n}=2i\int ds\,\partial_sW(\phi)=0$，符合预期。另一方面，若考虑半时空 $s\leq0$，则得到
\begin{equation}\tag{C.13}
  \int ds\,E^n_{2d}J_{2d,n}=2iW(\phi|),
\end{equation}
其中 $\phi|$ 是 $\phi$ 在边界 $s=0$ 上的值。当 $W(\phi|)$ 非零时，若希望构造与 $B$ 型扭变相容的边界条件，我们就不得不引入边界自由度；这有时称为 Warner 问题 [151]。该问题的标准解决办法是分解 $W$：引入边界 Fermi 多重态 $\Gamma^a_{1d}$，其 $E$ 型与 $J$ 型超势 $E^a_{1d}(\phi|)$、$J_{1d,a}(\phi|)$ 满足 $E^a_{1d}(\phi|)\cdot J_{1d,a}(\phi|)=-2iW(\phi|)$：
\begin{equation}\tag{C.14}
  S_{1d+2d}=S_{2d}+\frac{1}{4}\int dtd^2\theta\,\Gamma^a_{1d}\bar\Gamma_{1d,a}-\frac{1}{4}\int dtd\theta\,\Gamma^a_{1d}J_{1d,a}(\Phi|)\Big|_{\bar\theta=0}+\text{c.c.}
\end{equation}
我们看到总的 $E$ 项与 $J$ 项满足 $E^a_{1d}J_{1d,a}+\int ds\,E^n_{2d}J_{2d,n}=0$，因此组合作用量 $S_{1d+2d}$ 在由 $Q_B,\bar Q_B$ 生成的 1d $\mathcal{N}=2$ 超代数下不变。
% [Tutor Note: (C.13) 的边界项 $2iW(\phi|)$ 是 Warner 问题的核心：$B$ 型扭变把体超势变成边界上的全导数，若 $W(\phi|)\neq0$ 则边界条件被破坏。(C.14) 的因子化 $E^a_{1d}J_{1d,a}=-2iW$ 正是用边界 Fermi 多重态抵消这一边界项；它与 C.3 节的 $\delta_{\mathcal{E}}^2=W(\phi|)\mathrm{id}_{\mathcal{E}}$ 在 $W\to\frac{i}{2}W$ 的代换下是同一条件的不同记账方式。]

% ==== 原文 PDF 第 76 页 ====
\subsection{由 $B$ 膜得到矩阵分解}
现在我们已经有了定义 1/2-BPS $B$ 型边界条件（也称为 $B$ 膜）所需的经典数据，可以转向束缚在这类边界条件上的局域算子，更一般地，转向可用于在两个不同边界条件之间插值的局域算子，即 $B$ 膜范畴中的态射。习惯上只对整个 2d 手征多重态取 Neumann 边界条件，使边界值 $\phi|$ 不受约束，并对完整超势 $W(\phi)$ 的分解作研究。（下面我们作替换 $W\to\frac{i}{2}W$ 以简化表达式。）

如经典工作 [78] 中那样，借助态–算符对应，可以把两个边界条件交汇处的局域算子看作两侧各有一个边界条件的长条上的态；参见 Figure 3。我们不对该 Hilbert 空间作详细分析，而是从 Figure 3 左侧的视角以及一些代数来猜想结果。

% [背景补充: 矩阵分解 $(E,\delta_E)$ 是 Kontsevich 为 $B$ 型 Landau–Ginzburg 模型提出的 $B$ 膜代数模型：$\delta_E$ 满足 $\delta_E^2=W\mathrm{id}_E$，故它不是微分，而是一族在 $W$ 的临界点上才退化为微分的"扭曲"微分；矩阵分解的（DG-）范畴经局部化（形式化地商去"契约"对象）后给出 LG $B$ 膜的导出范畴。]
我们先从代数角度解释超势 $W$ 的分解 $(E^a,J_a)$。如 C.1 节那样，把边界费米子看作对边界值 $\phi|$ 编码一个模。然而现在映射 $\delta$ 不再作为微分起作用——它的平方是 $W(\phi|)$！特别地，若为模 $\mathcal{E}$ 选取一组基，则自同态 $\delta$ 是一个由 $\phi^n|$ 的多项式构成的矩阵，其平方等于 $2W$ 乘单位矩阵。我们把这样的 $\mathbb{C}[\phi^n|]$-模（或更一般地，凝聚层）$\mathcal{E}$ 连同 $R$ 荷／度为 $1$ 的自同态 $\delta_{\mathcal{E}}:\mathcal{E}\to\mathcal{E}$、满足 $\delta_{\mathcal{E}}^2=W(\phi|)\mathrm{id}_{\mathcal{E}}$ 的数据，称为 $W$ 的一个矩阵分解。
% [存疑: 同页正文先作 $W\to\frac{i}{2}W$ 的代换（第 76 页首段），故定义式的「平方为 $W(\phi|)$」与「平方等于 $2W$ 乘单位矩阵」（对应 (C.14) 的 $E^a_{1d}J_{1d,a}=-2iW$ 记账）宜视为同一条件的两种记账方式，而非笔误；此处照原文排印，未改。]

两个边界条件交汇处的局域算子必然与 $\phi^n|$ 的乘法对易：我们可以把 $\phi^n|$ 的插入从边界拉进体里、再在任意时刻 $t$ 拉回边界。因此，若两个边界条件分别由矩阵分解 $(\mathcal{E},\delta_{\mathcal{E}})$ 与 $(\mathcal{F},\delta_{\mathcal{F}})$ 编码，则它们交汇处的局域算子应是这些 $\mathbb{C}[\phi^n|]$-模之间的映射 $f:\mathcal{E}\to\mathcal{F}$；物理的局域算子将是那些 $\bar Q_B$-闭的模映射，模去 $\bar Q_B$-恰当的映射。虽然 $\delta_{\mathcal{E}}$ 与 $\delta_{\mathcal{F}}$ 不是微分，它们确实在矩阵分解之间的映射空间上诱导出一个微分 $\delta$（等同于 $\bar Q_B$ 的作用）：若 $f$ 的 $R$ 荷／度为 $r$，则度为 $r+1$ 的映射 $\delta f$ 由 $\delta f:=\delta_{\mathcal{F}}f-(-1)^r f\delta_{\mathcal{E}}$ 给出，从而
\begin{equation}\tag{C.15}
  \delta^2f=\delta_{\mathcal{F}}\delta f-(-1)^{r+1}\delta f\delta_{\mathcal{E}}=Wf-fW=0.
\end{equation}
% [Tutor Note: (C.15) 的要点是 $\delta^2 f=Wf-fW$：在模映射的层面 $f$ 两侧都作用着同一个超势 $W$，所以 $\delta$ 平方为零（$\delta$ 成为真正微分）当且仅当把 $f$ 视为 $W$ 的矩阵分解之间的映射；对单个矩阵分解，$\delta_{\mathcal{E}}^2=W\mathrm{id}$ 说明它只是一个" twisted" 微分。对偶地，矩阵分解范畴的 Hom 复形上确有 $\delta^2=0$，这正是 $B$ 膜范畴的 DG 结构来源。]
由此我们定义矩阵分解的（DG-）范畴 $MF(W)$ 为这样的范畴：其对象是 $W$ 的矩阵分解，从 $(\mathcal{E},\delta_{\mathcal{E}})$ 到 $(\mathcal{F},\delta_{\mathcal{F}})$ 的态射是 $\mathbb{C}[\phi^n|]$-模态射构成的（DG-）向量空间 $\operatorname{Hom}_{\mathbb{C}[\phi^n|]-\mathrm{mod}}(\mathcal{E},\mathcal{F})$，其微分为 $\delta f:=\delta_{\mathcal{F}}f-(-1)^r f\delta_{\mathcal{E}}$。

为使上述内容稍微具体些，下面简要描述一个简单边界条件上的局域算子空间，即一个简单矩阵分解中的自同态空间。考虑矩阵分解 $(\mathcal{E},\delta_{\mathcal{E}})$，其底层的 $\mathbb{C}[\phi^n|]$-模为 $\mathcal{E}=\mathbb{C}[\phi^n,\bar\gamma_a]$，自同态如上为 $\delta_{\mathcal{E}}=E^a\bar\gamma_a+J_a\partial_{\bar\gamma_a}$。$\mathbb{C}[\phi^n|]$-模映射 $\mathcal{E}\to\mathcal{E}$ 构成的空间（本身也看作 $\mathbb{C}[\phi^n|]$-模）由度为 $r_a$ 的费米子映射 $\partial_{\bar\gamma_a}$ 与度为 $-r_a$ 的费米子映射 $\bar\gamma_a$（即乘以 $\bar\gamma_a$）生成，也就是说，这个 $\mathbb{C}[\phi^n|]$-模映射空间就是 $\mathbb{C}[\phi^n|,\bar\gamma_a,\partial_{\bar\gamma_a}]$。该代数上的微分 $\delta$ 由与 $\delta_{\mathcal{E}}$ 作分次交换子诱导；在这些生成元上，它由下式给出
\begin{equation}\tag{C.16}
  \delta\phi^n|=0\qquad \delta\bar\gamma_a=J_a\qquad \delta\partial_{\bar\gamma_a}=E^a.
\end{equation}
把 $\phi^n|$ 替换为 $\phi^n$，这恰好就是 3.2.2 节给出的描述。
